% Draft sections for the OSNEM submission (addendum A1, docs/12-protocol-addendum.md). Every number is TODO until
% the post-hoc run of 99_final_eval writes results/posthoc_eval/tables/; cite the table file next to each number.
% \input this file from zero_shot_sexism_llms.tex after the main results, once the runs exist.

\section{Analyses added after the test labels were read}
\label{sec:posthoc}
The main results (Section~\ref{sec:results}) follow a protocol fixed on 2026-10-20 before any test label was read.
The analyses in this section were added afterwards. Their rules (metrics, thresholds, tie rules, word lists, prompt
variants and sample sizes) were written in a dated protocol addendum before any of them ran. We label each
analysis as \emph{post hoc} (a new view of the frozen predictions) or \emph{pre-registered} (new inference runs
whose rules were fixed before they ran). None of them changes a primary number. The test composition of CMSB
(adversarial modifications) was inspected after primary scoring.

\subsection{How certain are the comparisons?}
% Source: stats_display.csv, mde_by_task.csv, size_vs_f1_ci.csv, cross_task_tau.csv
We report simultaneous (Bonferroni over 14 comparisons) intervals for the difference to the best model next to the
per-comparison intervals of Table~\ref{tab:main}, the minimum detectable difference at 80\% power, and the
probability that each model is the best one. TODO: MDE per task; which rankings are stable.

\subsection{Why models miss sexism}
\label{sec:failure}
% Sources: failure_fnr_by_form.csv, failure_by_support.csv, failure_reason_shares.csv, failure_nonanswer_tokens.csv,
% failure_fpr_gender_terms.csv, failure_expert_override.csv, cmsb_contrast_set.csv, exist_dev_fnr_by_intention.csv
For every model, post and label we record whether the model said Yes, and, for every miss, a reason code: a
confident No (log-odds $\le -2$), a borderline No, or a non-answer (answer mass $< 0.5$, with the class of the
first reply token). We relate misses to the form of sexism (EDOS categories and vectors; explicit derogation and
threats vs implicit animosity and prejudiced discussion; CMSB source scales; EXIST source intention), to annotator
support, to the presence of slurs and gender terms, and to the CMSB adversarial pairs.
TODO: which forms are missed most; the share of confident vs borderline misses; whether false flags concentrate
on posts with gender terms or minority annotator support; contrast-set pair accuracy.

\subsection{Explainability}
\label{sec:xai}
% Sources: xai_faithfulness_by_cell.csv, xai_lexicon_effects.csv, xai_plausibility_cmsb.csv,
% xai_counterfactual_flips.csv, xai_self_explanation_codes.csv
We explain the Yes/No log-odds of the best models with SHAP and LIME on a stratified sample of true and false
positives and negatives. We test whether the attributions are faithful (ERASER comprehensiveness and sufficiency
against random words, with masking and with deletion) and plausible (overlap with the words the CMSB modification
changed). Two counterfactual rewrites probe what drives the decision on every scored post: swapping gendered words,
and removing slurs and profanity. We contrast these with the models' own one-sentence explanations of their
errors. Those explanations are model claims, not faithful accounts of the decision \citep{turpin2023}.
TODO: results; limits of each method.

\subsection{EXIST 2025 on the development split}
\label{sec:exist-dev}
% Sources: exist_dev_official_metrics.csv, exist_dev_strict.csv
The EXIST 2025 test gold is not public, and scoring English-only predictions against the bilingual gold is not
valid. We therefore score the zero-shot models on the 489 English development tweets with the official metrics
(ICM, ICM-Norm, F1; ICM-Soft for the soft setting; PyEvALL) and the official hierarchy for tasks 1.2 and 1.3.
No model is trained or tuned on these tweets, but they are a development split, the tweets are probably the EXIST
2023 corpus (TODO: confirm), and $n \approx 400$--$444$ cannot separate most models. We make no ranking claims here.
TODO: results with intervals.

\subsection{Robustness to prompt wording and option order}
% Sources: prompt5_*.csv, option_order_*.csv, answer_prefix_v2.csv, format_check.csv
TODO: spread over five frozen templates (mean $\pm$ SD, Kendall's $W$); permutation consistency and PriDe;
non-answer rate with an answer prefix; first-token vs free-text mismatch.

\subsection{Calibration, decision rules and annotator disagreement}
% Sources: threshold_free_binary.csv, label_free_decision_rules.csv, dev_tuned_rules_on_test.csv, soft_metrics.csv,
% risk_coverage.csv
TODO: AUROC/AUPRC; prior-matched and batch-calibrated decisions; dev-tuned thresholds; L1/Brier against annotator
shares; risk-coverage.

\subsection{Supervised reference}
% Sources: encoders_per_seed.csv, encoders_summary.csv
TODO: DeBERTa-v3-base and RoBERTa-base, three seeds each, mean $\pm$ SD; a reference point, not a claim of
superiority in either direction.

% ---------------------------------------------------------------- Elsevier front and back matter (TODO)
% Keywords (3-5): sexism detection; large language models; zero-shot classification; content moderation;
%   explainability   (TODO: confirm with co-authors)
% Highlights (3-5 bullets, <= 85 characters each): TODO after the results are final.
% Declaration of competing interest: TODO ("The authors declare no competing interests." if true).
% CRediT authorship contribution statement: TODO.
% Data availability: code, prompts, frozen configs, split ids and per-post model scores (no post text) on Zenodo
%   (TODO DOI); datasets from their original providers under their terms (TODO: check each licence).
% Declaration of generative AI and AI-assisted technologies in the writing process: TODO (required by Elsevier;
%   state which tools were used, for what, and that the authors take responsibility for the content).
